\documentclass{llncs}

\usepackage{standalone}
\input{preamble.tex}
\input{macros.tex}

\title{%
	Automatically Translating Proof Systems
	for SMT Solvers to the λΠ-calculus}
\author{Ciarán Dunne\inst{1,3} \and Guillaume Burel\inst{2,3}}
\institute{\scriptsize INRIA, ENS Paris-Saclay;
  \and ensIIE;
  \and SAMOVAR, Télécom SudParis, Institut Polytechnique de Paris, 91120 Palaiseau, France}

\begin{document}
\maketitle


\begin{abstract}
	\noindent
	Eunoia is a logical framework designed for
	specifying the proofs and proof systems of SMT solvers,
	namely \cvc{5}.
	%
	We present a translation from a core fragment of Eunoia
	to the \emph{$λΠ$-calculus modulo rewriting} as implemented
	by the LambdaPi proof assistant.
	%
	The translation is implemented by our tool \texttt{eo2lp},
	which we use for generating LambdaPi encodings of
	%
	(a)~a large fragment of the
	\textit{Cooperating Proof Calculus} (CPC),
	the Eunoia signature defining \cvc{5}'s proof system, and
	%
	(b)~proofs produced by \cvc{5} on \texttt{unsat}
	problems from various fragments of SMT-LIB.
\end{abstract}

\documentclass[class=llncs, crop=false]{standalone}
\input{preamble.tex}
\input{macros.tex}

\begin{document}
\section{Introduction}\label{sec:intro}
%
The area of automated reasoning known as
\emph{satisfiability modulo theories} (SMT)
focuses on the development of tools for
checking the satisfiability of first-order
formulas over combinations of mathematical
theories~\cite{Barrett2021}.
%
Such tools (known as SMT \emph{solvers}) are
increasingly used as back-ends for various tasks:
hardware~\cite{DBLP:conf/fmcad/MattareiMBDHH18} and
software verification~\cite{barnett06boogie,riant23debugging,kosmatov16frama-c},
%
model checking \cite{DBLP:conf/cav/ChampionMST16},
cyber-security \cite{Backes2018,rungta22billion},
and to improve automation in proof
assistants~\cite{DBLP:conf/cpp/ArmandFGKTW11,DBLP:journals/jar/BlanchetteBP13}.
%
Because these applications demand a high level
of confidence in the results produced by the solver,
many solvers are capable of generating
\emph{proof certificates}
that witness the unsatisfiability of formulas
reported as such by the solver.
%

% Towards the aim of standardizing and benchmarking solvers,
% the \emph{SMT library initiative} oversees the development of
The \emph{SMT-LIB standard}~\cite{Barrett2015-standard}
specifies the syntax and semantics of a
common language used for interacting with solvers.
%%
Proof generation is however not covered by the standard,
leading to the development of various proof formats,
namely LFSC~\cite{DBLP:journals/fmsd/StumpORHT13},
Alethe~\cite{schurrAletheGenericSMT2021},
and recently Eunoia~\cite{EthosUser_manualmdMain}.
%
Eunoia is both a proof output language and a logical framework
for specifying the constants, side-condition programs,
and inference rules of an SMT solver (e.g., \cvc{5}).
The \emph{Cooperating Proof Calculus} (CPC) is the
Eunoia signature encoding \cvc{5}'s proof system,
and Eunoia proof scripts may be checked by the C++ tool
Ethos~\cite{EthosUser_manualmdMain}.
%
%

The $\lambda\Pi$-calculus modulo rewriting~\cite{CousineauDowek}
and its implementations in the tools Dedukti~\cite{assaf2023dedukti}
and LambdaPi~\cite{hondet:hal-02981561} provide
trustworthy proof checking with an emphasis on
interoperability of proof systems.
%
This paper presents a translation from Eunoia to LambdaPi
and an OCaml tool $\texttt{eo2lp}$ that implements the
translation.
%
We used our tool to translate \emph{CPC-MINI}, a fragment
of CPC covering 14~SMT-LIB logic fragments
(\autoref{fig:pipeline}), and to further translate and
typecheck over 1{,}000~proof scripts produced by running
\cvc{5} on $\texttt{unsat}$ problems from those fragments.
% ---------------------------------------------------------
\begin{boxfigure}[t!]{fig:pipeline}
	{{Overview of the translation pipeline.
				\cvc{5} produces Eunoia proof scripts using rules
				from the CPC signature, which may be checked by Ethos.
				\texttt{eo2lp} translates both the signature and the proofs
				to LambdaPi, which independently verifies them.}}
	%
	\centering
	\alttext{Pipeline diagram with three vertical groups.
		Left: SMT-LIB fragments (QF\_UF, QF\_LIA, UFLIA, ALIA, QF\_NIA, UF, \ldots).
		Middle: Eunoia artefacts---a CPC signature (theories, rules, programs)
		and a Eunoia proof script (.prf.eo).
		Right: LambdaPi artefacts---a translated CPC signature plus Prelude,
		and a LambdaPi proof file (.prf.lp).
		Horizontal arrows: cvc5 maps SMT-LIB fragments to Eunoia proof scripts;
		eo2lp maps both the Eunoia signature and proof script to LambdaPi.
		Downward arrows: ethos checks the Eunoia bundle, lambdapi check verifies
		the LambdaPi bundle, each producing a tick-or-cross result.}{%
		\input{fig-pipeline.tex}}
\end{boxfigure}
% ---------------------------------------------------------

% \subsubsection{Contribution.}
% Practically,
% \texttt{eo2lp} (\autoref{sec:results}) allows
% using LambdaPi as an independent proof checker for
% \cvc{5}: it translates all 22 modules of CPC-MINI
% (359 rules, 96 programs, 57 constants) and verifies
% 97\% of 1{,}036 proof scripts across 14 SMT-LIB
% fragments.
% %
% Conceptually, the translation gives an interpretation
% of Eunoia in the $λΠ$-calculus in the absence of a
% model-theoretic semantics like those given in the
% SMT-LIB standard.
% %
% Eunoia's (dependent) types are interpreted as terms of a
% type universe~$\Set$, inference rules as axioms,
% and side-conditions as rewrite rules.

\subsubsection{Related Work.}
The reduction of SMT proof checking to type
checking originates with
LFSC~\cite{DBLP:journals/fmsd/StumpORHT13};
Eunoia inherits this approach and extends it with
higher-order features similar to those of the
speculative SMT-LIB~3
proposal~\cite{barrettSMTLIB3Bringing,smt3-proposal,barbosa17hosmt}.
%
Brown~\cite{Brown2021} compares the prospects of interpreting
the proposed type system within type-theoretic and set-theoretic
frameworks; our encoding realizes the former via the $λΠ$-calculus
modulo rewriting.

The Dedukti/LambdaPi ecosystem is an interoperability hub for
proof assistants (Rocq~\cite{boespflug12coqine},
Isabelle~\cite{isabelle_dedukti},
HOL~Light~\cite{assaf15translating},
Matita~\cite{krajono},
Lean~\cite{vaishnav25lean4less})
and automated provers
(iProverModulo~\cite{burel20first},
ArchSAT~\cite{bury:tel-02612985},
Vampire~\cite{komel2025case},
TPTP~\cite{DBLP:conf/flairs/SutcliffeBB25}).
%
Coltellacci et al.~\cite{coltellacciReconstructionSMTProofs2024a,coltellacci25frocos}
reconstruct cvc5 proofs in LambdaPi via hand-written soundness lemmas for
each rule of the Alethe~\cite{schurrAletheGenericSMT2021} proof format.
%
Similarly, \textsc{lean-smt}~\cite{mohamed25leansmt}
reconstructs Eunoia proofs in Lean 4 using lemmas
that encode a subset of CPC's rules.
%
Our work instead targets Eunoia at the framework
level to allow encoding arbitrary
Eunoia signatures and proof scripts
in a type-theoretic framework.


\subsubsection*{Outline}
%
\autoref{sec:eunoia} introduces the fragment of Eunoia
targeted by our translation and its elaboration step;
\autoref{sec:lambdapi} defines the translation into the
$λΠ$-calculus modulo rewriting uniformly across rules,
programs, and proof scripts;
\autoref{sec:results} describes \texttt{eo2lp} and
reports the benchmark results.

%
% SMT proof reconstruction is also used within
% proof assistants:
% Isabelle/HOL replays Alethe
% proofs~\cite{DBLP:conf/cade/SchurrFD21,lachnitt25isabelle},
% SMTCoq~\cite{ekici17smtcoq} integrates solvers
% into Rocq, and IsaRare~\cite{lachnitt24isarare}
% verifies CPC rewrite rules in Isabelle/HOL.


% ---------------------------------------------------------
% Placed here so it floats to the top of page 2
\begin{boxfigure}[t!]{fig:eo-syntax}
	{Eunoia syntax: terms, attributes, built-in symbols, commands, proof commands.}
	\alttext{Five stacked grammar tables.
		First: terms---symbols, literals, applications, or binder applications;
		literals as int(n), rat(n,m), or str(s); parameters and program cases.
		Second: parameter attributes (:implicit, :list) and constant attributes
		(right/left-associative with optional nil terminator, chainable, pairwise,
		binder, arg-list).
		Third: built-in symbols grouped by family---core (eo::eq, eo::requires,
		eo::quote, eo::var), boolean (eo::and, eo::or, eo::not, eo::xor),
		arithmetic (eo::add, eo::mul, eo::neg, eo::zdiv), and
		list (eo::nil, eo::cons, eo::list\_concat, eo::list\_len).
		Fourth: standard commands---declare-const, declare-parameterized-const,
		declare-consts, define (macro definition), program (with signature and
		cases), declare-rule (with optional premises/premise-list, args,
		requires, and conclusion attributes), include.
		Fifth: proof commands---assume, assume-push, step (with rule, optional
		premises, optional args), and step-pop.}{%
		\vspace{4pt}\centerline{$\eoTermSyntax$}
		\boxfigurehline
		\centerline{$\eoAttrSyntax$}
		\boxfigurehline
		\centerline{$\eoBuiltinFig$}
		\boxfigurehline
		\centerline{$\eoCommandSyntax$}
		\boxfigurehline
		\centerline{$\eoPrfCommandSyntax$}\vspace{4pt}}
\end{boxfigure}
% ---------------------------------------------------------

\end{document}


\documentclass[class=llncs, crop=false]{standalone}

\input{preamble.tex}
\input{macros.tex}

\begin{document}
\section{Eunoia}\label{sec:eunoia}
% fig:eo-syntax is in intro.tex so it floats to the top of page 2
% ---------------------------------------------------------
This section presents the fragment of Eunoia targeted by our
translation, together with the elaboration step that resolves
$n$-ary applications to normal form.
%
Eunoia currently has no published formal semantics:
its meaning is defined operationally by Ethos.
%
The presentation here was informed by the
Ethos user manual~\cite{EthosUser_manualmdMain},
the Ethos source, and through discussions with the
\cvc{5}/Eunoia developers.
%
% We accordingly do not give a typing relation or
% algorithm for Eunoia, nor a model-theoretic semantics
% analogous to the one in SMT-LIB~2.6~\cite{Barrett2015-standard}.


% ---------------------------------------------------------
\subsection{Syntax, Commands, and Signatures}\label{sec:eo-sig}
% ---------------------------------------------------------
%
\Autoref{fig:eo-syntax} defines the concrete syntax for Eunoia,
where the metavariable $s$ ranges over a set of \emph{symbols}.
%
A term is either a symbol $s$, a literal $ℓ$
(integer, rational, or string), an \emph{application}
$\eapp{s}{\vec t}$, or a \emph{binder application}
$\eapp{s}{\epar{\vec ρ}\ t}$.
%
A \emph{parameter} $\epar{s\ t\ \maybe{ν}}$ declares a symbol
$s$ of type $t$ with an optional \emph{parameter attribute} $ν$.
Similarly, a symbol may have a \emph{constant attribute} $α$
that specifies how its $n$-ary applications are elaborated
to binary form.
%
Eunoia provides a fixed set of \emph{built-in symbols}
$β$ that operate on the syntactic structure of terms;
their role is explained alongside the constructs that
use them (rule declarations in \autoref{sec:eo-rules}
and elaboration in \autoref{sec:eo-elab}).
%
Commands are divided into \emph{standard commands}
and \emph{proof commands}.

% ------------------------------------------------------
\begin{table}[t!]
	\centering
	\caption{Declarations produced by each Eunoia command.
		Parts in $\maybe{\cdot}$ are present only when the
		corresponding attribute is supplied. Each program case
		$c_i = \ecase{l_i}{r_i}$ contributes one rewrite rule.}
	\label{tab:cmd-decls}
	\small
	\begin{tabular}{@{}>{$}p{0.56\linewidth}<{$}@{\ }>{$}p{0.40\linewidth}<{$}@{}}
		\hline
		\textnormal{\textbf{Command}} & \textnormal{\textbf{Declarations}} \\
		\hline
		\epar{\kw{declare-const}\ s\ t\ \maybe{α}}
		                              & \{s : t\}\,\maybe{∪\,\{s ∷ α\}} \\
		\epar{\kw{declare-parameterized-const}\ s\ \epar{\vec ρ}\ t\ \maybe{α}}
		                              & \{s\prn{\vec ρ} : t\}\,\maybe{∪\,\{s ∷ α\}} \\
		\epar{\kw{declare-consts}\ ℓ\ t}
		                              & \textnormal{binds the literal class~$ℓ$ to type~$t$} \\
		\epar{\kw{define}\ s\ \epar{\vec ρ}\ t\ \maybe{\attr{type}\ t'}}
		                              & \{s\prn{\vec ρ} ↪ t\}\,\maybe{∪\,\{s\prn{\vec ρ} : t'\}} \\
		\epar{\kw{program}\ s\ \epar{\vec ρ}\ \attr{signature}\ \epar{\vec t}\ t'\ \epar{c_1 \ldots c_m}}
		                              & \{s\prn{\vec ρ} : \earr{\vec t\ t'}\} ∪ \{l_i ↪ r_i \mid c_i = \ecase{l_i}{r_i}\} \\
		\epar{\kw{declare-rule}\ s\ \epar{\vec ρ}\ \ldots}
		                              & \{s\prn{\vec ρ} : τ^★\}\textnormal{\ (\autoref{sec:eo-rules})} \\
		\epar{\kw{include}\ μ}
		                              & \textnormal{declarations of module~$μ$} \\
		\epar{\kw{assume}\ s\ φ}\textnormal{\,/\,}\epar{\kw{assume-push}\ s\ φ}
		                              & \{s : \eprf{φ}\} \\
		\epar{\kw{step}\ s\ φ\ \attr{rule}\ r\ \maybe{\attr{args}\ \epar{\vec t}}\ \maybe{\attr{premises}\ \epar{\vec p}}}
		                              & \{s : \eprf{φ},\ s ↪ \eapp{r}{\vec t\ \vec p}\} \\
		\hline
	\end{tabular}
\end{table}
%
A \emph{signature} $Σ$ is a set of \emph{declarations},
where each declaration is either a \emph{typing} $s\prn{\vec ρ} : t$,
an \emph{attribution} $s\prn{\vec ρ} ∷ α$,
or a \emph{rewrite rule} $ℓ ↪ r$ with parameters $\vec \rho$.
%
\autoref{tab:cmd-decls} maps each command to
the set of declarations it contributes.
%
Any sequence of standard commands $\vec δ$ thus
determines a signature: the union of the
declarations contributed by each command.
%
A \emph{proof script} $Υ = \prn{\vec δ ⨾ \vec π}$
pairs a sequence of standard commands $\vec δ$
(the \emph{preamble}) with a sequence of proof
commands $\vec π$ (the \emph{body}); $Υ$ likewise
determines a signature, which we call a \emph{proof}.
%

\subsection{Rule Declarations and Proofs}\label{sec:eo-rules}
%
Two of the core built-in symbols play a role specific
to rule declarations.
%
The guard $\eoreqs{t_1}{t_2}{t_3}$ evaluates to $t_3$
only when $t_1$ and $t_2$ are syntactically identical,
and is otherwise stuck; it lets a rule make its
conclusion conditional on a side-condition.
%
The marker $\eoquote{t}$ flags $t$ as a \emph{quoted
	argument} whose value may appear in the return type
of the rule, giving Eunoia's mechanism for
object-level dependent types.
%
%
Given an inference rule declaration
$\epar{\kw{declare-rule}\ s\ \epar{\vec ρ}\ \ldots\ \conclusion φ}$,
the derivation of its contributed type declaration
$s\prn{\vec\rho} : \tau^\star$ is given thus:
%
\begin{enumerate}
	\item Given
	      $\requires{\ecase{l_1}{r_1}\ldots\ecase{l_k}{r_k}}$,
	      define $φ^★$ below. Otherwise, $φ^★ ≔ φ$.
	      $$φ^★ ≔ \eoreqs{l_1}{r_1}{
			      \epar{
				      \,\ldots\,\eoreqs{l_k}{r_k}{φ}}
		      }
	      $$

	\item If the declaration provides no premises,
	      let $τ^★ ≔ \eprf{φ^★}$.
	      Otherwise, define $τ^★$ thus:
	      \begin{enumerate}
		      \item Given
		            $\attr{premises}\ \epar{\plur ψ m}$,
		            let:
		            $$τ^★ ≔ \earr{\eprf{ψ_1}\,\ldots\,\eprf{ψ_m}\ \eprf{φ^★}}$$
		            %
		      \item Given $\premiselist ψ f$ for
		            some associative symbol $f$, redefine:
		            $$φ^★ ≔\eoreqs{\eapp{\eo{is\_list}}{f\ ψ}}{\ev{true}}{φ^★}$$
		            and let $τ^★ ≔ \earr{\eprf{ψ}\ \eprf{φ^★}}$.
	      \end{enumerate}

	\item If the declaration provides
	      $\attr{args}\ \epar{\plur t n}$,
	      then the type contributed for $s$ is:
	      $$s\prn{\vec ρ}
		      : \earr{\eoquote{t_1}\,\ldots\,\eoquote{t_n}\ τ^★}
	      $$
	      Otherwise, $s\prn{\vec ρ}: τ^★$.
\end{enumerate}

\begin{example}\label{ex:eo-rules}
	Consider four rule declarations from
	\texttt{cpc/rules/Booleans.eo}.
	%
	\begin{eobox}
		\lstinputlisting[style=eo, aboveskip=0pt, belowskip=0pt]{listings/rule-modus-ponens.eo}%
		\listsep
		\lstinputlisting[style=eo, aboveskip=0pt, belowskip=0pt]{listings/rule-cnf-implies-pos.eo}%
		\listsep
		\lstinputlisting[style=eo, aboveskip=0pt, belowskip=0pt]{listings/rule-reordering.eo}%
		\listsep
		\lstinputlisting[style=eo, aboveskip=0pt, belowskip=0pt]{listings/rule-and-intro.eo}%
	\end{eobox}
	%
	\noindent
	The $\ev{modus\_ponens}$ rule has only $\attr{premises}$ and
	$\attr{conclusion}$, so item~2a gives:
	$$
		\ev{modus\_ponens}(\vec ρ)  :
		{\color{eoPunct}{\texttt{(}}}\ev{->}\
		\eprf{\ev{F1}}\
		\eprf{\eapp{\ev{=>}}{\ev{F1}\ \ev{F2}}}\
		\eprf{\ev{F2}}
		{\color{eoPunct}{\texttt{)}}}
	$$
	%
	The $\ev{cnf\_implies\_pos}$ rule has
	$\attr{args}\ \epar{\eapp{\ev{=>}}{\ev{F1}\ \ev{F2}}}$;
	item~3 makes this an
	$\eo{quote}$-typed parameter that pattern-matches at
	the point of use, where $φ$ is the term given by
	its $\attr{conclusion}$ attribute:
	$$
		\ev{cnf\_implies\_pos}(\vec ρ)  :
		{\color{eoPunct}{\texttt{(}}}\ev{->}\
		\eoquote{\eapp{\ev{=>}}{\ev{F1}\ \ev{F2}}}\
		\eprf{φ}
		{\color{eoPunct}{\texttt{)}}}
	$$
	%
	The $\ev{reordering}$ rule combines $\attr{args}$
	with $\attr{requires}$: item~1 builds the guarded
	conclusion
	$φ^★ ≔ \eoreqs{\eapp{\eo{list\_minclude}}{\ev{or}\ \ev{C1}\ \ev{C2}}}{\ev{true}}{\ev{C2}}$,
	item~2a adds the premise $\ev{C1}$, and item~3 wraps
	the $\eo{quote}$-typed argument $\ev{C2}$:
	$$
		\ev{reordering}(\vec ρ)  :
		{\color{eoPunct}{\texttt{(}}}\ev{->}\
		\eoquote{\ev{C2}}\
		\eprf{\ev{C1}}\
		\eprf{φ^★}
		{\color{eoPunct}{\texttt{)}}}
	$$
	%
	The $\ev{and\_intro}$ rule has
	$\premiselist{\ev F}{\ev{and}}$, so item~2b guards
	the conclusion with
	$φ^★ ≔ \eoreqs{\eapp{\eo{is\_list}}{\ev{and}\ \ev F}}{\ev{true}}{\ev F}$
	and takes the premise list as a single $\eprf{\ev F}$
	argument:
	$$
		\ev{and\_intro}(\vec ρ)  :
		{\color{eoPunct}{\texttt{(}}}\ev{->}\
		\eprf{\ev F}\
		\eprf{φ^★}
		{\color{eoPunct}{\texttt{)}}}
	$$
\end{example}


% ---------------------------------------------------------
\subsection{Elaboration}\label{sec:eo-elab}
% ---------------------------------------------------------
%
Unlike standard LF-style frameworks, Eunoia supports
$n$-ary applications of binary symbols, with the
unrolling determined by a \emph{constant attribute}
$α$ (\autoref{fig:eo-syntax}) assigned to the symbol
at its declaration.
%
After elaboration, every application of a declared
constant $s$ uses the built-in higher-order
application operator $\mtt{\_}$, so that
$\eapp{s}{t_1\ t_2}$ becomes $\eho{\eho{s}{t_1}}{t_2}$;
only built-in and program symbols are applied
directly.
%
For example, given $f$ with attribute $\rcn{t_\nil}$,
the \emph{$f$-list} with elements $\plur t n$ is the
nested application
$\eho{\eho{f}{t_1}}{\eho{\eho{f}{\ldots}}{\eho{\eho{f}{t_n}}{t_\nil}}}$.
%
The built-in symbols $\eo{nil}$, $\eo{cons}$, and
$\eo{list\_concat}$ construct and combine such
$f$-lists from within programs and rule conclusions.

A second elaboration strategy targets binder
applications.  Eunoia provides a built-in
\emph{heterogeneous list} type $\eo{List}$, with
constructors $\eo{List::nil}$ and $\eo{List::cons}$,
whose entries may have different types, and a built-in
symbol $\eo{var}$ that constructs variable-like terms
$\eapp{\eo{var}}{\estr{x}\ t}$.
%
The attribute $\bd{t_\cons}$ specifies that a binder
application $\eapp{s}{\epar{\vec ρ}\ t}$ is elaborated
so that each parameter $\epar{x_i\ t_i}$ in $\vec ρ$
becomes an $\eo{var}$ term, and these terms are
gathered with the list constructor $t_\cons$ (typically
$\eo{List::cons}$, with terminator $\eo{List::nil}$).

The $\mbf{elab}$ operator below captures both
strategies uniformly.  It is parameterized by a
signature $Σ$ and a local parameter context $\vec ρ$;
we write $\msf{progs}(Σ)$ for the set of symbols
introduced by $\kw{program}$ commands in $Σ$.

\newcommand{\body}{\text{body}}
\newcommand{\vars}{\text{vars}}
\begin{definition}\label{def:elab}
	For any signature $Σ$ and parameter context $\vec ρ$,
	let $\elab{Σ,\vec ρ}$ be the least operator such that:
	\begin{enumerate}
		\item For any symbol $s$, $\elab{Σ,\vec ρ}[s] = s$.
		      %
		\item For a binder application
		      $\eapp{s}{\epar{\vec v}\ t}$ where
		      $\prn{s ∷ \bd t_\cons} ∈ Σ$, each parameter
		      $\epar{x_i\ t_i} ∈ \vec v$ is reified as
		      $X_i ≔ \epar{\eo{var}\ \estr{x_i}\ t_i}$,
		      and the $X_i$ are gathered with the list
		      constructor $t_\cons$:
		      $$ \elab{Σ,\vec ρ}[\eapp{s}{\epar{\vec v}\ t}]
			      =
			      \eapp{s}{t_\vars\ t_\body}
			      \quad\text{where}\quad
			      \begin{array}{l@{\ {≔}\ }l}
				      t_\vars &
				      \elab{Σ, \vec ρ}[\eapp{t_\cons}{\plur X n}]
				      \\
				      t_\body &
				      \elab{Σ, \vec ρ}
				      [\mbf{subst}_{\prn{\vec x, \vec X}}[t]]
			      \end{array}
		      $$

		      %
		\item For an application $\eapp s {\vec t}$,
		      let $t_i' ≔ \elab{Σ, \vec ρ}[t_i]$ in:
		      \begin{enumerate}
			      \item If $s$ is a built-in or
			            $s ∈ \msf{progs}(Σ)$, then
			            $\elab{Σ, \vec ρ}[\eapp{s}{\vec t}] =
				            \eapp{s}{\vec t'}$.
			      \item If $s ∈ \vec ρ$, then
			            $\elab{Σ, \vec ρ}[\eapp{s}{\vec t}] =
				            \eho{\ldots\eho{\eho{s}{{t_1'}}}{\ldots}}{{t_n'}}$.

			      \item If $\prn{s ∷ α} ∈ Σ$ for an
			            associative $α$, the result is
			            one of:
			            $$
				            \begin{scases}
					            \foldr(G, t_\nil, \vec t')
					            & \tsm{$\prn{α = \rcn{t_\nil}}$}
					            \\
					            \foldl(G, t_\nil, \vec t')
					            & \tsm{$\prn{α = \lcn{t_\nil}}$}
					            \\
					            \foldr(G, t'_n, t'_1\ldots t'_{n-1})
					            & \tsm{$\prn{α = \rc}$}
					            \\
					            \foldl(G, t'_1, t'_2\ldots t'_{n})
					            & \tsm{$\prn{α = \lc}$}
				            \end{scases}
			            $$
			            where the combining function
			            $G(x,y)$ is $\econcat{s}{x}{y}$ if
			            $\prn{x ∷ \attr{list}} ∈ \vec ρ$
			            and $\eho{\eho{s}{x}}{y}$
			            otherwise.

			      \item If $\prn{s ∷ α} ∈ Σ$ for a
			            non-associative $α$, then
			            $\elab{Σ,\vec ρ}[\eapp{s}{\vec t}] =
				            \elab{Σ,\vec ρ}[\eapp{t_\cons}{\vec u}]$
			            where:
			            $$
				            \vec u ≔
				            \begin{scases}
					            \eapp{s}{t_1 t_2}\,\eapp{s}{t_2 t_3}
					            \ldots\,\eapp{s}{t_{n-1} t_n}
					            &
					            \tsm{$\prn{α = \ch {t_\cons}}$}
					            \\
					            \eapp{s}{t_i t_j} \text{ for all } i < j
					            &
					            \tsm{$\prn{α = \pw {t_\cons}}$}
					            \\
					            \vec t
					            &
					            \tsm{$\prn{α = \al {t_\cons}}$}
				            \end{scases}
			            $$

		      \end{enumerate}
	\end{enumerate}
\end{definition}

\begin{example}\label{ex:elab}
	The constants $\ev{or}$, $\ev{=}$, and $\ev{forall}$
	carry attributes that drive each of the elaboration
	cases:
	%
	\begin{eobox}
		\lstinputlisting[style=eo, aboveskip=0pt, belowskip=0pt]{listings/decls-or-eq.eo}%
		\listsep
		\lstinputlisting[style=eo, aboveskip=0pt, belowskip=0pt]{listings/decl-forall.eo}%
	\end{eobox}
	%
	\noindent
	Since $\ev{or}$ has attribute $\rcn\false$, item~3c
	applies with $G(x,y) = \eho{\eho{\ev{or}}{x}}{y}$ and
	the right fold gives:
	$$ \elab{Σ, \vec ρ}[\eapp{\ev{or}}{\ev x\ \ev y\ \ev z}]
		=
		\eho{\eho{\ev{or}}{\ev x}}{
			\eho{\eho{\ev{or}}{\ev y}}{
				\eho{\eho{\ev{or}}{\ev z}}{\false}}}
	$$
	%
	Since $\ev{=}$ has attribute $\ch{\ev{and}}$, item~3d
	first chains the equalities and then recurses
	through $\ev{and}$'s $\rcn\false$ attribute:
	$$
		\elab{Σ, \vec ρ}[\eapp{\ev{=}}{\ev a\ \ev b\ \ev c\ \ev d}]
		=
		\elab{Σ, \vec ρ}[\eapp{\ev{and}}{
				\eapp{\ev{=}}{\ev a\ \ev b}\
				\eapp{\ev{=}}{\ev b\ \ev c}\
				\eapp{\ev{=}}{\ev c\ \ev d}}]
	$$
	%
	Since $\ev{forall}$ has attribute
	$\bd{\eo{List::cons}}$, item~2 reifies each parameter
	as an $\eo{var}$ term inside an $\eo{List}$ and
	elaborates the body under the substitution
	$\sigma ≔ \prn{\ev{x} ↦ X}$, where
	$X ≔ \epar{\eo{var}\ \estr{x}\ \ev{Int}}$:
	$$
		\begin{array}{r@{\ {=}\ }l}
			\elab{Σ, \vec ρ}[\eapp{\ev{forall}}{\epar{\epar{\ev{x}\ \ev{Int}}}\ \ev{F}}]
			&
			\eapp{\ev{forall}}{t_\vars\ t_\body} \\[1mm]
			t_\vars &
			\eapp{\eo{List::cons}}{X\ \eo{List::nil}} \\[1mm]
			t_\body &
			\elab{Σ, \vec ρ}[σ\ev{F}]
		\end{array}
	$$
\end{example}

\end{document}



\documentclass[class=llncs, crop=false]{standalone}
\input{preamble.tex}
\input{macros.tex}

\begin{document}
%
\section{Encoding Eunoia in the λΠ-calculus}\label{sec:lambdapi}
\begin{boxfigure}[t!]{fig:lpTermSyntax}
	{Syntax and typing rules for the $λΠ$-calculus modulo rewriting.}
	\alttext{Top: term grammar for the lambda-Pi calculus modulo rewriting.
		Universes mu range over TYPE and KIND.
		Terms t range over variables, constants, universes, applications,
		lambda abstractions, and dependent products.
		Bottom: natural-deduction typing rules---empty context is well-formed
		(wf0); extending with (x:t) requires Gamma deriving t at some universe
		(wf+); variable rule reads (x:t) from the context (var);
		constant rule reads (kappa:t) from the signature with the constant's
		type sortable in Gamma (con); TYPE has type KIND (univ);
		dependent product, abstraction, application, and conversion modulo
		beta and rewriting (prod, fun, app, conv).}{%
		\vspace{4pt}\centerline{$\lpTermSyntax$}
		\boxfigurehline
		\centerline{$\lpTypingRules$}\vspace{4pt}}
\end{boxfigure}
%
\Autoref{fig:lpTermSyntax} provides syntax and typing rules
for the \emph{$λΠ$-calculus modulo rewriting}.
%
Each term is either a \emph{variable} $x$,
a \emph{constant} $κ$,
a \emph{universe} from $\{ \TYPE, \KIND \}$,
an \emph{application} of two terms $\prn{\app{t}{t'}}$,
or an \emph{abstraction}
$\prn{\binddot{\mscr{B}}{x : t}{t'}}$
where $\mscr{B}$ is a \emph{binder} from $\{ λ, Π \}$.
%
Terms are identified up to $α$-conversion
(i.e., renaming of bound variables), and we assume
appropriate definitions for \emph{substitution} $\prn{t[x ↦ t']}$
and \emph{$β$-reduction} $\prn{t ⇝_{β} t'}$.
%
Hereinafter, we may write $\prn{t_1 → t_2}$ for the term
$\prn{\ppii{(x : t_1)}{t_2}}$, where $x$ does not occur
free in $t_2$.

% --------------------------------------------------------%
A \emph{typing} is an expression of the form $\prn{t : t'}$,
and a \emph{context} is a list of typings each of the
form $\prn{x : t}$.
%
% Hereinafter, given a context $Γ$ and a typing $\prn{x : t}$,
% let $\prn{Γ, \prn{x : t}}$ denote the extension $Γ ∪ \{x : t\}$.
%
A \emph{rewrite rule} is an expression of the form
$\prn{t ↪ t'}$, where $t$ has the form
$\prn{\app{\prn{\app{κ}{t_1}}\ \ldots}{t_n}}$.
%
A \emph{signature} is a list of typings
$\prn{κ : t}$ (with $t$ closed) and rewrite rules.
%
For any signature $Σ$, let $R_{Σ}$ be the smallest
binary relation such that:
%
\begin{enumerate}
	\item for any rule $\prn{ℓ ↪ r} ∈ Σ$ and substitution $σ$,
	      $(ℓσ, rσ) ∈ R_{Σ}$, and
	\item $R_{Σ}$ is \emph{congruent} under application
	      and abstraction.
\end{enumerate}
%
%, abstraction and substitution.  $R$
Then, \emph{equality modulo rewriting} $\prn{{≡_{βΣ}}}$ is
defined as the least equivalence relation containing
$R_{Σ}$ and the $β$-reduction relation.
%
The rules in \autoref{fig:lpTermSyntax}
provide a (mutually inductive) definition of a
\emph{well-formedness} relation $\prn{\msf{wf}_Σ}$
on contexts and a \emph{typing relation}
$\prn{⊢_{Σ}}$ between contexts and typings,
where $\prn{Γ ⊢_{Σ} e : t}$ may be read as
``$e$ has type $t$ with respect to signature $Σ$ and context $Γ$''.

% -------------------------------------------------------- %
\subsection{Translating Eunoia Terms}
%
The typing rules of the $λΠ$-calculus disallow
abstractions on types (i.e., there are no expressions
of type $\TYPE → t$ in any context). Eunoia's types are
therefore encoded as $λΠ$-terms of type $\Set$, using
the Tarski-style universe of LambdaPi's standard
library: an injective symbol
$\msf{τ} : \Set → \TYPE$ lifts a $\Set$-term to a
$λΠ$-type, and the (infix) function-space symbol
$\prn{⤳} : \Set → \Set → \Set$ satisfies
$\msf{τ}\,\prn{A ⤳ B} ↪ \msf{τ}\,A → \msf{τ}\,B$.
%
The dependent types arising from $\eo{quote}$ arguments
are not expressible with $\prn{⤳}$ alone, so we
additionally introduce a dependent function space
$\prn{⤳^{\msf d}} :
	\pii{T : \Set}{\prn{{\msf{τ}\,T → \Set}} → \Set}$
with the rule
$\msf{τ}\,\prn{T ⤳^{\msf d} F} ↪
	\pii{x : \msf{τ}\,T}{\msf{τ}\,\prn{\app F x}}$;
the non-dependent case is recovered by taking
$F ≔ \lam{\_}{B}$.
%
We keep $\prn{⤳}$ as a first-class operator rather than
defining it from $\prn{⤳^{\msf d}}$ because programs
need a matchable function-space head: pattern variables
cannot bind under a $Π$, but they can under $\prn{⤳}$.

The translation $\tm{⋅}$ on Eunoia terms below produces
$λΠ$-terms in the abstract calculus of
\autoref{fig:lpTermSyntax} and assumes all terms have
already been processed by the $\mbf{elab}$ operator of
\autoref{sec:eo-elab}.
%
Because some Eunoia symbols are not valid LambdaPi
identifiers, the translation on symbols uses LambdaPi
syntax for escaping: $\bar s ≔ \esc s$ iff $s$
contains forbidden characters, and $\bar s ≔ s$
otherwise.
%
\begin{definition}\label{def:tm-ty}
	%
	Let $\tm{⋅}$ be the least operator such that:

	\begin{enumerate}
		\item For any symbol $s$ or literal $ℓ$, the following hold:
		      $$\tm{s} =
			      \begin{scases}
				      \lpo{Type}   & \text{if $s = \eotype$,} \\
				      \bar s & \text{otherwise.}
			      \end{scases}
			      %
			      \quad
			      %
			      \tm{ℓ} =
			      \begin{scases}
				      \ev{of\_Z}\ \tm{n}
				      &
				      \text{if $ℓ = \msf{int}(n)$,}
				      \\
				      \ev{of\_Q}\ \tm{n}\ \tm{m}
				      &
				      \text{if $ℓ = \msf{rat}(n,m)$,}
				      \\
				      \estr{s} & \text{if $ℓ = \msf{str}(s)$.}
			      \end{scases}
		      $$
		      The symbols $\ev{of\_Z}$ and $\ev{of\_Q}$ are
		      defined in our prelude (\autoref{sec:results}) and embed
		      the integer and rational literals of LambdaPi's
		      standard library into the carrier type of Eunoia's
		      arithmetic sorts.
		\item For arrow types, the following hold,
		      where $\epar{x\ T} ∈ \vec ρ$:
		      $$
			      \begin{array}[t]{r@{\ {=}\ }l}
				      \tm{\earr{\eoquote{x}\ {\vec t}}}
				                    &
				      \tm{T} ⤳^{\msf d} \lam{x}{\tm{\earr{\vec t}}}
				      \\[1mm]
				      \tm{\earr{t'\ {\vec t}}}
				                    &
				      \tm{t'} ⤳ \tm{\earr{\vec t}}
				      \\[1mm]
				      \tm{\earr{t}} & \tm t
			      \end{array}
		      $$
		\item For binary and $n$-ary applications,
		      the following hold:
		      $$
			      \begin{array}[t]{l@{\ {=}\ }l}
				      \tm{\eapp{\mtt{\_}}{t_1\ t_2}}
				       & \lpo{app}\ \tm{t_1}\ \tm{t_2}
				      \\[2mm]
				      \tm{\eprf{t}}
				       & \lpo{Proof}\ {\tm {t}}
				      \\[2mm]
				      \tm{\eapp{s}{\vec t}}
				       &
				      \bar s\ \tm{t_1}\ \ldots\ \tm{t_n}
			      \end{array}
		      $$
		\item After elaboration, every other binder
		      application is unfolded by item~2 of
		      \autoref{def:elab}; only $\eo{define}$
		      survives because it is a let-binder rather
		      than a binding form. It is translated using
		      LambdaPi's standard $\msf{let}$ binder:
		      $$\tm{\eapp{\eo{define}}{
					      \epar{\vec v}
					      \ t'
				      }
			      }
			      = \prn{
				      \msf{let}\ \vec {w}
				      \ \msf{in}\ \tm{t'}
			      }$$
		      where $w_i = \prn{{\bar x}_i ≔ \tm{t_i}}$
		      for each $v_i = \epar{x_i\ t_i}$.
	\end{enumerate}
	%
\end{definition}

% -------------------------------------------------------- %
\subsection{The LambdaPi Proof Assistant}
%
\begin{boxfigure}[t!]{fig:lpSyntax}
	{Syntax for a fragment of the LambdaPi proof assistant.}
	\alttext{Two grammar tables for LambdaPi concrete syntax.
		Top: modifiers (constant, sequential); parameters (explicit (s:t) and
		implicit [s:t]); patterns theta (pattern variables \$x or pattern
		applications); rewrite rules (s applied to patterns rewriting to theta).
		Bottom: commands---symbol declaration with optional modifier, parameters,
		type, and optional definition; rule declaration listing rewrite rules with
		'with' separators; assertion 'assert |- t : t';
		and 'require open' for module imports.}{%
		\vspace{4pt}\centerline{$\lpPatternSyntax$}
		\boxfigurehline
		\centerline{$\lpConcreteSyntax$}\vspace{4pt}}
\end{boxfigure}
%
\Autoref{fig:lpSyntax} defines the fragment of LambdaPi
required for encoding Eunoia. It is sugar over the
calculus of \autoref{fig:lpTermSyntax}: implicit binders
$\imp{x : t}$ and placeholders $\_$ are elaborated to
$Π$-binders and metavariables during type inference, and
each declaration extends~$Σ$.

A \emph{symbol declaration}
$\maybe{m}\,\kw{symbol}\ s\ \vec ρ : t\,\maybe{≔ t'}$
introduces $s$ with the typing
$\prn{\xsym s : \pii{\vec ρ}{t}}$, where parameters
$\vec ρ$ are either \emph{explicit} $\xpl{x : t}$ or
\emph{implicit} $\imp{x : t}$.
We write $\xsym s$ (with an @) for the
\emph{fully-explicit} form; an unprefixed~$s$ inserts a
placeholder for each leading implicit, to be filled by
higher-order unification.
The optional $\md{constant}$ modifier forbids rewrite
rules with $s$ at the head; $\md{sequential}$ matches
rewrite rules with head $s$ in their textual order; and
a definition $\prn{≔ t'}$ adds the rewrite rule
$\prn{s ↪ \lam{\vec ρ}{t'}}$.
%
A \emph{rewrite rule declaration} $\kw{rule}\ θ ↪ θ'$
adds the rule $\prn{θ ↪ θ'}$ to the signature, where
$θ$ and $θ'$ are \emph{patterns}---either a
\emph{pattern variable} $\ptrn{x}$ (written \$x) or a
\emph{pattern application} $\prn{s\,\any{θ}}$.
The rule is valid when $θ$ is a pattern application and
every pattern variable in $θ'$ also occurs in $θ$.
LambdaPi additionally verifies \emph{subject reduction}
at declaration time: for every typed substitution of
the pattern variables, both sides must have the same
type. An \emph{assertion} $\kw{assert}\ ⊢\ t : t'$
checks $t : t'$ without extending the signature.

% ----------------------------------------------------------
\subsection{Translating Eunoia Commands and Proofs}
% ----------------------------------------------------------
We extend the term translation $\tm{⋅}$
(\autoref{def:tm-ty}) to parameter contexts and to
both standard and proof commands.
%
\Autoref{tab:encoding} summarises the full
Eunoia-to-LambdaPi encoding; the remainder of this
section formalises the operators on contexts
(\autoref{def:ctx}), standard commands
(\autoref{def:cmd}), and proof commands
(\autoref{def:cmd-prf}).

\begin{table}[t!]
	\centering
	\caption{Eunoia-to-LambdaPi encoding.
		Symbols prefixed with $\msf{eo.}$ are defined in
		our hand-written prelude (\autoref{sec:results});
		$⤳$ and $⤳^{\msf d}$ are the function-space
		symbols introduced in \autoref{sec:lambdapi}.}
	\label{tab:encoding}
	\small
	\begin{tabular}{@{}p{0.40\linewidth}@{\ \ }p{0.52\linewidth}@{}}
		\hline
		\textbf{Eunoia feature} & \textbf{LambdaPi encoding} \\
		\hline
		Type universe $\eotype$
		& term $T : \Set$, carrier $\msf{τ}\,T$ \\
		Function type $\earr{A\ B}$
		& $A ⤳ B$, with $\msf{τ}\,(A ⤳ B) ↪ \msf{τ}\,A → \msf{τ}\,B$ \\
		Dependent type from $\eoquote{x}$
		& $A ⤳^{\msf d} \lam{x}{B}$ (matchable head, unlike $Π$) \\
		Parameter attribute $\attr{implicit}$
		& implicit binder $\imp{x : \msf{τ}\,T}$ \\
		Macro $\kw{define}$
		& $\kw{symbol}\ s\,\vec ρ ≔ t$ \\
		Higher-order application $\eho{t_1}{t_2}$
		& $\lpo{app}\ t_1\ t_2$ \\
		$n$-ary $\eapp{f}{\vec t}$ with attribute
		& nested binary application; nil terminator via rewrite rule \\
		Binder variables ($\eo{var}$)
		& reified into the list type $\lpo{List}$ \\
		Program ($\kw{program}$)
		& $\md{sequential}\ \kw{symbol}$ + one rewrite rule per case (SR-checked) \\
		Inference rule ($\kw{declare-rule}$)
		& $\md{constant}\ \kw{symbol}$ + auxiliary symbol with quote-pattern rewrite when $\attr{args}$ present \\
		Built-in arithmetic
		& prelude rules delegating to $\mathbb Z, \mathbb Q$ \\
		Conclusion guard $\eoreqs{l}{r}{ψ}$
		& $ψ$ applies only when $l ≡ r$ syntactically \\
		Proof command $\kw{assume}$
		& $\md{constant}\ \kw{symbol}$ of type $\lpo{Proof}\,φ$ \\
		Proof command $\kw{step}$
		& defined symbol + $\kw{assert}$ \\
		\hline
	\end{tabular}
\end{table}


% ----------------------------------------------------------
\subsubsection{Translating Commands.}
% ----------------------------------------------------------
Parameter contexts and two auxiliary notations
support the command translation that follows.

\begin{definition}\label{def:ctx}
	Given $\vec ρ = \epar{x_1\ t_1\ \maybe{ν_1}}
		\ldots \epar{x_n\ t_n\ \maybe{ν_n}}$, define:
	$$
		\ctx{\vec ρ}\ ≔\
		ρ_1^★\, \ldots\, ρ_n^★
		\quad\text{where}\quad
		ρ_i^★ =
		\begin{scases}
			\imp{\bar x_i : \El{\tm{t_i}}}
			& \text{if $ν_i = \attr{implicit}$,} \\
			\xpl{\bar x_i : \El{\tm{t_i}}}
			& \text{otherwise.}
		\end{scases}
	$$
\end{definition}
%
We further write $\tm{t}^{\ptrn{}}$ for the result of
applying $\tm{⋅}$ to $t$ and replacing each free
variable $x$ with the pattern variable $\ptrn{x}$;
and $\mbf{impl}(\vec ρ, t)$ for the sublist of
$\ctx{\vec ρ}$ containing each $x_i ∈ \vec ρ$ that
occurs free in $t$, marked as implicit.
%
\begin{definition}\label{def:cmd}
	%
	Let $δ$ be a standard Eunoia command. Then $\cmd{δ}$
	is the list of LambdaPi commands described as follows:
	%
	\begin{enumerate}\itemsep3pt
		\item If $δ$ is a constant declaration,
		      $δ ⊢ s(\vec ρ) : t$, the head of
		      $\cmd δ$ is
		      $\kw{symbol}\ \bar s\ \ctx{\vec ρ}
			      : \El{\tm t}\mtt{;}$.
		      If $δ ⊢ s ∷ α$ and $α$ is
		      $\rcn{t_\nil}$ or $\lcn{t_\nil}$, the
		      following rewrite rule is also produced,
		      where $\lpo{nil}$ returns the nil-terminator
		      of an associative symbol and $\lpo{as}$
		      performs type ascription:
		      $$
			      \kw{rule}\ \prn{\lpo{nil}\ \bar s\ \ptrn T}
			      ↪ \prn{\lpo{as}\ \tm{t_\nil}\ \ptrn T}
		      $$

		\item If $δ$ is a macro definition,
		      $δ ⊢ s(\vec ρ) ≔ t$, then
		      $\cmd{δ}$ is
		      $\kw{symbol}\ \bar s\ \ctx{\vec ρ} :
			      \El{\tm{t'}} ≔ \tm{t}$
		      (with the typing omitted if no $t'$
		      is supplied).

		\item If $δ$ is a program declaration with
		      $δ ⊢ s(\vec ρ) : \earr{\plur t n\, t'}$
		      and cases $\ecase{l_1}{r_1}\ldots\ecase{l_m}{r_m}$,
		      then $\cmd{δ}$ declares $\bar s$,
		      where $T ≔ \El{\tm{ \earr{\plur t n\, t'}}}$:
		      %
		      $$\begin{array}[t]{l}
				      \md{sequential}\ \kw{symbol}\ \bar s\
				      \vec ρ^★ : T
				      \mtt{;}
				      \\[1mm]
				      \begin{array}[t]{@{}l@{\ }l@{\ {↪}\ }l}
					      \kw{rule} & \tm{l_1}^{\ptrn{}} & \tm{r_1}^{\ptrn{}}
					      \\ \cdots \\
					      \kw{with} & \tm{l_m}^{\ptrn{}} & \tm{r_m}^{\ptrn{}}\mtt{;}
				      \end{array}
				      % \quad (j = 1 \ldots m)
			      \end{array}$$
		      %
		      where $\vec ρ^★ ≔ \mbf{impl}(\vec ρ, T)$ binds
		      all of the free parameters in $T$ as implicits.

		\item If $δ$ is a rule declaration with
		      $\attr{args}\ \epar{\plur t n}$ of type
		      $\plur τ n$, then
		      $δ ⊢ s(\vec ρ) :
			      \earr{
				      \eoquote{t_1}\,\ldots\,\eoquote{t_n}
				      \ τ^★}$,
		      where $τ^★$ encodes the premises and
		      conclusion (\autoref{sec:eo-sig}).
		      An auxiliary symbol $s^★$ eliminates the
		      $\lpo{quote}$ types by rewriting to
		      $\tm{τ^★}$ on terms matching the argument
		      patterns; $s$ then takes the arguments
		      explicitly and uses $s^★$ for its return type:
		      $$
			      \begin{array}[t]{l}
				      \kw{symbol}\ \bar{s}^★\
				      \vec σ : \tm{τ_1} → \ldots → \tm{τ_n} → \TYPE \mtt{;}
				      \\[1mm]
				      \kw{rule}\
				      \bar s^★\,\tm{t_1}^{\ptrn{}}\,\ldots\,\tm{t_n}^{\ptrn{}}
				      ↪ \tm{τ^★}^{\ptrn{}}\mtt{;}
				      \\[1mm]
				      \md{constant}\ \kw{symbol}\ \bar s\
				      (α_1 : \tm{τ_1}) \ldots (α_n : \tm{τ_n}) :
				      \prn{\app{\app{{\app{\bar s^★}{α_1}}}{\ldots}}{α_n}}
			      \end{array}
		      $$
	\end{enumerate}
\end{definition}

\begin{example}\label{ex:cmd}
	We illustrate $\cmd{δ}$ on the four declarations
	from \autoref{ex:eo-rules}:
	%
	\begin{lpbox}
		\lstinputlisting[style=lp, aboveskip=0pt, belowskip=0pt]{listings/cmd-constants.lp}%
		\listsep
		\lstinputlisting[style=lp, aboveskip=0pt, belowskip=0pt]{listings/cmd-modus-ponens.lp}%
		\listsep
		\lstinputlisting[style=lp, aboveskip=0pt, belowskip=0pt]{listings/cmd-rule-aux.lp}%
	\end{lpbox}
	%
	\noindent
	The polymorphic equality $\ev{=}$ and the
	associative $\ev{or}$ translate by case~1, the latter
	also producing a nil-terminator rewrite rule.
	$\ev{modus\_ponens}$ has no $\attr{args}$, so case~1
	yields its derived type directly.
	$\ev{cnf\_implies\_pos}$ has the compound argument
	$\attr{args}\epar{\eapp{\ev{=>}}{\ev{F1}\ \ev{F2}}}$,
	triggering case~4: an auxiliary symbol matches the
	argument pattern and rewrites to the translated
	conclusion under $\lpo{Proof}$.
\end{example}

% ----------------------------------------------------------
\subsubsection{Translating Proofs.}
% ----------------------------------------------------------

\begin{definition}\label{def:cmd-prf}
	%
	Let $π$ be a proof command. Then $\cmd{π}$
	is the list of LambdaPi commands such that:
	%
	\begin{enumerate}\itemsep6pt
		\item If $π$ is an assumption,
		      then $π ⊢ s : \eapp{\mtt{Proof}}{φ}$:
		      %
		      $$\md{constant}\ \kw{symbol}\ \bar s :
			      \lpo{Proof}\ \tm{φ}\mtt{;}$$

		\item If $π$ is a step,
		      then $π ⊢ s : \eapp{\mtt{Proof}}{φ}$
		      and $π ⊢ s ≔ \eapp{r}{\plur t m\ \plur p n}$:
		      %
		      \begin{align*}
			       & \kw{symbol}\ \bar s
			      ≔ \bar r\ \tm{t_1}\ \ldots\ \tm{t_m}\
			      \bar p_1\ \ldots\ \bar p_n\mtt{;}
			      \\
			       & \kw{assert}\ ⊢\ \bar s :
			      \lpo{Proof}\ \tm{φ}\mtt{;}
		      \end{align*}
		      %
		      The body applies rule $\bar r$ to its
		      arguments then premises.
		      %
		      The type is checked separately via
		      $\kw{assert}$, which uses LambdaPi's
		      type inference to fill implicit arguments.
	\end{enumerate}
\end{definition}

\begin{example}\label{ex:proof}
	A \texttt{QF\_UF} proof script (top) deriving
	$\false$ from $\eapp{\ev =}{\ev{x0}\ \ev{x0}}$
	via $\ev{eq{-}refl}$ and $\ev{eq\_resolve}$ (CPC
	uses both \texttt{-} and \texttt{\_} in rule names),
	and its LambdaPi translation (bottom) generated by
	\autoref{def:cmd} and \autoref{def:cmd-prf}:
	%
	\begin{eobox}
		\lstinputlisting[style=eo, aboveskip=0pt, belowskip=0pt]{listings/proof.eo}%
		\listsep
		\lstinputlisting[style=lp, aboveskip=0pt, belowskip=0pt]{listings/proof.lp}%
	\end{eobox}
	%
	The preamble's $\ev{U}$, $\ev{x0}$ translate by
	case~1 and the $\kw{define}$s by case~2 of
	\autoref{def:cmd}; the assumption $\ev{@p1}$
	becomes a $\md{constant}$ symbol of type
	$\lpo{Proof}\ \tm{φ}$ (case~1 of
	\autoref{def:cmd-prf}); each step applies its
	rule to arguments then premises, with $\kw{assert}$
	verifying the conclusion via type inference.
\end{example}

\end{document}


\documentclass[class=llncs, crop=false]{standalone}

\input{preamble.tex}
\input{macros.tex}

\begin{document}
\section{Implementation and Results}\label{sec:results}

Our tool
\texttt{eo2lp}\footnote{\url{https://github.com/deducteam/eo2lp}}
is an OCaml program
(${\sim}5{,}500$~lines) that implements
\autoref{sec:lambdapi}, with a Menhir-generated
parser~\cite{menhir} and a back-end built against
the LambdaPi OCaml API.
%
A hand-written \texttt{Prelude.lp} module
(${\sim}335$~lines) bootstraps the type embedding:
it declares $\Set$, $\msf{τ}$, the function space~$⤳$,
the higher-order application symbol $\lpo{app}$,
Boolean logic, and arithmetic builtins as rewrite
rules over LambdaPi's integers/rationals
(e.g.\ $\eo{add}\,(\ev{of\_Z}\,i)\,(\ev{of\_Z}\,j) ↪
	\ev{of\_Z}\,(i + j)$).
The prelude also declares the built-in heterogeneous
list type $\lpo{List} : \Set$ (the encoding of Eunoia's
$\eo{List}$), together with its constructors
$\lpo{List::nil}$ and $\lpo{List::cons}$. These hold
the $\eapp{\eo{var}}{\estr{x_i}\ T_i}$ entries produced
by $\mbf{elab}$ on binder applications
(\autoref{sec:eo-elab}, item~2).
Closed expressions reduce concretely, so CPC's
\texttt{\$run\_evaluate} program---which recursively
delegates to the $\eo{add}, \eo{mul}, \eo{and},
	\eo{not},\ldots$ builtins---normalises ground terms
like $\eapp{\ev{+}}{2\ 3}$ to $5$.

\subsection{CPC-MINI}
The Cooperating Proof Calculus (CPC) is the
Eunoia signature that formalizes \cvc{5}'s proof
system.
%
The full CPC comprises 35~core modules
(${\sim}11{,}000$~lines of Eunoia) plus
16~expert modules for additional theories
(separation logic, bags, finite fields,
floating points, transcendentals).
%
To validate our approach on a manageable portion,
we define
\emph{CPC-MINI},\footnote{\url{https://github.com/ciaran-matthew-dunne/cpc-mini}}
a fork of the upstream CPC included as a git
submodule of \texttt{eo2lp}.
%
It is a modified subset of
22~core modules (${\sim}5{,}000$~lines)
covering logic, integer arithmetic, arrays,
sets, quantifiers, and uninterpreted functions
--- 6~theory modules (57~constants),
8~rule modules (353~rules, 50~programs), and
7~program modules (38~programs, 36~macros),
plus a root module \texttt{Cpc} that adds
6~further rules (e.g.\ \texttt{evaluate},
\texttt{trust}) for 359~rules in total.

CPC-MINI applies three classes of change to the
upstream CPC.
\emph{(i)~Theory pruning:} the expert modules and
unsupported core theories (bitvectors, strings,
datatypes, reals) are removed, since their primitives
are not yet in our prelude.
\emph{(ii)~Type-parameter tightening:} a few CPC
declarations leave their numeric parameter type
unconstrained so that the same rule applies at both
\texttt{Int} and \texttt{Real}; we instantiate these
to \texttt{Int} where required, which suffices for
our integer-arithmetic benchmarks.
\emph{(iii)~Bug fixes:} several upstream typos and
mismatched type parameters are patched; we intend to
contribute these upstream. For instance, the
upstream \texttt{\$mk\_cong\_rhs} program in
\texttt{Uf.eo} typed its arguments with the codomain
rather than the domain of its function
parameter---a mismatch Ethos missed because it only
typechecks program clauses at application time.
%
All 22~modules pass
\texttt{lambdapi check} with full subject-reduction
verification; the entire CPC-MINI signature is
translated in ${\sim}100$\,ms and type-checked in
${\sim}160$\,ms---this is one-off work, independent
of the proofs to be verified later.

\subsection{Proof Benchmarks}
%
To evaluate the translation on proofs,
we sampled up to 100~unsatisfiable problems from
SMT-LIB for each of the 14~fragments covered by
CPC-MINI and obtained Eunoia proof scripts from
\cvc{5}, yielding 1{,}036~proofs.
%
Each proof is translated by \texttt{eo2lp} and
the result is checked by \texttt{lambdapi check}
with subject-reduction verification,
under a 5-second timeout.
%
\Autoref{tab:results} summarizes the results:
1{,}002 of 1{,}036~proofs (97\%) are
successfully translated and verified,
with an average checking time of 0.47\,s.

\input{results-table.tex}

The seven equality and array fragments achieve a
near-perfect pass rate (582 of 583), with the single
failure being a timeout. The seven linear-arithmetic
fragments pass 420 of 453 proofs (93\%):
QF\_UFIDL, UFLIA, and LIA reach 98--100\%; the
remaining failures are 18 timeouts in QF\_LIA and
QF\_IDL where large proofs exceed the 5-second
limit, and 15 type-checking errors in QF\_UFLIA
caused by CPC rules whose type signatures mix
\texttt{Real} and \texttt{Int}---most prominently
\texttt{arith-int-geq-tighten}, which declares a
\texttt{Real} parameter but is applied here in an
integer-only context.
Such mixed-arithmetic rules are excluded by
CPC-MINI's type-parameter tightening (change~(ii)
above); the failure is therefore a known
\emph{type coverage} gap at the
\texttt{Real}/\texttt{Int} boundary, not a semantic
mismatch with Ethos. Crucially, \emph{zero} failures
are caused by a semantic disagreement: every proof
that Ethos accepts and we translate is accepted by
LambdaPi up to coverage and time.

% To probe scaling, we ran the translator on a larger
% sample of 439 \texttt{QF\_UF} proofs (up to
% ${\sim}5{,}000$ lines each): all 439 check
% successfully, with median check time 0.49\,s,
% $p_{95}$ 1.79\,s, and a single outlier of 32.9\,s on
% a proof generating long rewrite chains. The encoding's
% overhead over pure type-checking is one $\msf{τ}$
% indirection per type plus rewriting for program
% evaluation; pure type-checking is linear, but the
% rewrite-based program evaluation needed for
% side-conditions and arithmetic can blow up on
% unusually long chains. Outliers in the main
% benchmark are dominated by such chains.

\end{document}


\section{Conclusion and Future Work}\label{sec:conclu}

We have presented an automatic translation from
Eunoia to the $\lambda\Pi$-calculus modulo rewriting,
unlike approaches that hand-write soundness lemmas per
rule~\cite{coltellacciReconstructionSMTProofs2024a,mohamed25leansmt}.
Translated rules are axioms, but their well-typedness
and subject reduction are verified by LambdaPi;
\texttt{eo2lp} translates all 22 modules of CPC-MINI
and verifies 97\% of 1{,}036 proof scripts.
Relative to Ethos, our pipeline is slower
(0.47\,s vs.\ ${<}0.1$\,s) and supports a subset of CPC,
but adds two guarantees Ethos does not---rule-signature
well-typedness and subject reduction of side-condition
programs---the latter surfaced upstream CPC bugs.
Future work includes extending CPC-MINI to the full CPC,
proving CPC rules sound directly in
$λΠ$~\cite{lachnitt24isarare}, and using the typed
interpretation of \autoref{sec:lambdapi} as a setting
for the higher-order extensions proposed for
SMT-LIB~3~\cite{barrettSMTLIB3Bringing,smt3-proposal}.

\subsubsection*{Acknowledgements.}
%
Research reported in this document was supported by an
Amazon Research Award,~2022, and by the CARMA associate
team, whose funding for collaboration visits helped the
progress of this work.
%
The authors thank Haniel Barbosa,
Hans-Jörg Schurr, Andrew Reynolds, and Yoni Zohar
for discussions on the design of Eunoia and CPC, and
on earlier iterations of this work.
%
We also thank Frédéric Blanqui and Théo Winterhalter for
their comments on the paper, and the rest of Deducteam
for their support.
%
This paper is dedicated to the memory of Gilles Dowek,
who initiated this research project in particular and
more generally the ecosystem of Dedukti.

\subsubsection*{Disclosure of Interests.}
The authors have no competing interests to declare
that are relevant to the content of this article.

\printbibliography

\end{document}


% \hrule
% \vspace{2mm}

% 	\hrule
% $$\begin{array}{r@{\ }l@{\quad}l}
% 		δ &
% 		\begin{array}[t]{c@{\ }l}
% 			{⦂} & \prn{\decl{x}{\prn{\vec ρ}}{\prn{{:}\ t}}{\prn{{∷}\ α}}}
% 			∣ \prn{\defn{x}{\paren{\vec ρ}}{\paren{{≔}\ t}}{\paren{{:}\ t'}}}    \\
% 			{∣} & \paren{\prog{x}
% 			{\paren{\vec ρ}}
% 			{\paren{{:}\ {\paren{t_1 \ldots t_n} → t'}}}
% 			{\paren{r_1 \ldots r_m}}
% 			}                                                                      \\
% 			{∣} & \paren{\irule{x}
% 				{\paren{\vec ρ}}
% 				{\paren{φ_1 \ldots φ_n}}
% 				{\paren{t_1 \ldots t_m}}
% 				{\paren{ψ}}
% 			}
% 		\end{array}
% 		  & \mcomment{(declarations)}
% 	\end{array}
% $$
